\section{The exact unforced background}
\label{short-bg:section}

We use the real axisymmetric class without swirl.  Set
\[
 p=11/4,\quad c=\nu^{-2},\quad R=\epsilon^{-1},\quad
 \epsilon=e^{-d/\nu},\quad P=2\epsilon c,
 \qquad 0<d_-\le d\le d_+.
\]
Powers of $\nu^{-1},c$ and fixed polynomial cutoff lengths are denoted
by $\mathcal P_\nu$. Constants may depend on $d_\pm$ and on a fixed
derivative budget. The phase measure is normalized by $d\tau/P$.
Use the spaces $Y_p,D_p$ of Lemma~\ref{bg:common-body-graph}, and write
$A_0=L$ and $X_p=D_p\cap\{\operatorname{div}w=0\}$ with the same graph
norm. Let $v_c$ be the train of Proposition~\ref{sct:train} and
$V_0=U_{\rm raw}$ its divergence-free lift from Proposition~\ref{sct:lift},
with residual $F_R$. In particular,
\eqref{sct:residual-bound} gives
$\|F_R\|_{C_\tau W^{1,p}}\le\mathcal P_\nu\epsilon^{3/11}$.

For $m>8$, write $X_m=L^2(\mathbb R^2,\langle x\rangle^{2m})$.
The frozen growing pair and complementary bound are given by
Corollary~\ref{cor:short-seed-dilation}. The train estimates give
\begin{align}
 &\sup_{s\ge0}\sum_{j=0}^3\|D^jv_c(s)\|_\infty\le C,
 \label{short-bg:train-uniform}\\
 &\int_{\Sigma/\nu}^\infty
       (\|Dv_c(s)\|_\infty+\|D^2v_c(s)\|_\infty)\,ds
 \le {C\over\nu(1+\Sigma)}+C\nu^3.
 \label{short-bg:train-late}
\end{align}
The weighted coefficient and polynomial spatial $L^{22/3}$ bounds follow
from Lemma~\ref{sct:heat}. In normalized cylinder coordinates set
\begin{equation}
 H_p=L^p(s^{-9/8}\,ds\,dq/c\,d\eta),\qquad
 \|\sqrt{s}\,b\|_{Y_p}=(\pi R^2)^{1/p}\|b\|_{H_p},
 \quad
 T_0F(s)=\int_0^s e^{\nu(s-t)\Delta}\mathbb P_0 F(t){dt\over t}.
 \label{short-bg:measure}
\end{equation}
The global free graph and compact-output comparison are
\eqref{bg:profile-graph} and Proposition~\ref{bg:lp-free-comparison};
their inputs have arbitrary age support.

\subsection{One mixed construction for the finite age interval}

We first verify the common actual-cylinder comparison. On the centered
cell let $X_{m,c}=L^2(C_c,(1+q^2+\eta^2)^m)$ and let $J_c,R_c$ be
zero extension and restriction. Extend the free cylinder evolution by
the identity on the exterior summand of $X_m$ and the interaction by
zero there. The vorticity equation is
\begin{equation}
 \omega_s=(\nu\Delta-v_c(s)\cdot\nabla+A_c(s))\omega,
 \qquad A_c(s)\omega=-(K_c\omega)\cdot\nabla E_c(s).
 \label{short-bg:actual-cylinder}
\end{equation}
Here $E_c=\operatorname{curl}v_c$ and $K_c$ includes the horizontal
mean mode. On $s\le\Sigma/\nu$, the train comparison gives, in any
fixed finite weighted derivative budget,
\[
 \|v_c-V_c^p\|+\|E_c-E_c^p\|
 \le\delta_c:=\mathcal P_\nu
                  e^{C_\Sigma/\nu-b_\Sigma c^2}.
\]
The central radial Gaussian transports the radial weight exactly;
the velocities induced by the other periodic copies, the region near the
cell boundary, and the displayed error contribute at most
$C(\delta_c+e^{-b_\Sigma c^2})$ to its logarithmic derivative.
The weight is continuous across the cell boundary and its normal derivative has
a negative jump. Its distributional Laplacian is therefore bounded
above by $C_m$ times the weight. Weighted integration by parts gives
\begin{equation}
 \|S_{\nu,c}(t,s)\|_{X_{m,c}\to X_{m,c}}
 \le\exp[(C_m\nu+C\delta_c+Ce^{-b_\Sigma c^2})(t-s)].
 \label{short-bg:free-cylinder}
\end{equation}
One may smooth the weight, preserving uniform comparability, to justify
this calculation. On a fixed block, Duhamel for a smooth compactly
supported datum against the frozen radial transport gives convergence:
the transported support remains a fixed disk, the viscosity error is
$\nu\Delta$ applied only to that smooth datum, and the coefficient
change is $O(\nu L)+o(1)$ there. Density and
\eqref{short-bg:free-cylinder} give strong convergence on $X_m$,
uniformly for block lengths in a compact interval. The frozen radial
transport is unitary in this radial weighted space. Multiplying by its
inverse reduces the free convergence to $W_\nu\to I$ strongly with
$\|W_\nu\|\le1+o(1)$. The identity for
$\|W_\nu^*f-f\|^2$ then gives strong adjoint convergence too. This
argument makes no operator-norm assertion about the Laplacian.

The cylinder kernel is uniformly bounded away from its common local
$1/|x-y|$ singularity. On $|x|,|y|\le c/4$, paired periodic images
give
$|K_c(x-y)-K_{\mathbb R^2}(x-y)|
\le Cc^{-2}(1+|x|+|y|)$.
Multiplication by the Gaussian observation coefficient $\nabla E$
and weighted Cauchy--Schwarz control this part by $Cc^{-2}\|f\|_{X_m}$.
For $|y|>c/4$, the same Cauchy--Schwarz estimate gives a factor
$c^{1-m}$, with an additional harmless source-moment power if needed;
for $|x|>c/4$ the coefficient is Gaussian-small. The coefficient error
contributes $C\delta_c$. Thus
\begin{equation}
 \|J_cA_c(\sigma/\nu)R_c-A_{\rm pl,\sigma}\|_{X_m\to X_m}
 \le C_\Sigma(c^{-2}+\delta_c),\qquad 0\le\sigma\le\Sigma.
 \label{short-bg:interaction-cylinder}
\end{equation}
The near kernel gains one local derivative, by the Fourier multiplier
bound for $DK$, and the coefficient is localized. Truncate source and
observation, use local Rellich compactness, and then the same tail
estimates: every interaction is compact and their family is
collectively compact. The artificial exterior is invariant and has no
growing component. Restriction removes it. To implement zero
circulation in the comparison one may adjoin a radial mass-one
neutral vector to the plane zero-circulation space; radial Euler
linearization annihilates it. Cylinder curl sources and periodized
angular-two entrance vectors have zero circulation, so their actual
solutions remain in the desired subspace.

Fix $\Sigma$ and put $T=\Sigma/\nu$. Choose
$0<a<a'<b'<b<\gamma/(1+\Sigma)^2$. On a bounded block of length $L$,
the common compact-Duhamel lemma gives
\[
 U(t+L,t)=W_{\nu,t}e^{L L_{\nu t}}+o_{\nu\to0}(1)
 \quad\hbox{in operator norm},
\]
where $\|W_{\nu,t}\|\le1+o(1)$ and $W_{\nu,t},W_{\nu,t}^*\to I$
strongly. The error is uniform for block lengths in $[0,L]$.
Indeed the interaction Duhamel term converges in norm; its limit is
compact, so multiplying it by $W_{\nu,t}-I$ gives a norm-small error.
Let $E_t,F_t$ be the frozen growing and complementary spaces. Finite
rank and strong adjoint convergence make both off-diagonal blocks
small. For $L$ fixed sufficiently large, each full block has matrix
\[
 \begin{pmatrix}a_n&b_n\\c_n&d_n\end{pmatrix},\qquad
 \|a_n^{-1}\|\le A,\quad \|d_n\|\le D,\quad
 \|b_n\|+\|c_n\|\le\eta_\nu\longrightarrow0,
 \qquad AD<1.
\]
Here $A,D$ can be chosen below $e^{-b'L},e^{a'L}$ with strict
margins. Frozen projections use the periodized plane direct vectors,
which have zero circulation, and the restricted plane dual vectors,
with their finite pairing matrix normalized. Weighted-dual convergence
suffices; no far-field expansion of
the dual vector is used.

For completeness, construct the invariant subspaces using
\begin{equation}
 \xi_{n+1}=(c_n+d_n\xi_n)(a_n+b_n\xi_n)^{-1},\qquad
 \zeta_n=(a_n-\zeta_{n+1}c_n)^{-1}(\zeta_{n+1}d_n-b_n),
 \quad \xi_0=0,\quad\zeta_N=0.
 \label{short-bg:graphs}
\end{equation}
On a ball of radius $r=C\eta_\nu$, these maps have Lipschitz constant
at most $A(D+\eta_\nu r)/(1-A\eta_\nu r)<1$. They preserve the ball
if $r(1-AD)\ge A\eta_\nu(1+r^2)$. Hence the graphs are complementary,
their projections are uniformly bounded, and the products of their
coordinate maps have bounds $e^{a'(t-s)}$ forward on the center and
$e^{-b'(s-t)}$ backward on the growing space. Coordinate equivalence
is used only at the two endpoints of a product. An appended block of
length at most $L$ has uniformly bounded growing inverse and forward
map and changes only the prefactor.

At intermediate times in a block $[t_n,t_{n+1}]$, put
$J_t=U(t,t_n)|_{E_n^u}$ and
$F_t=\Pi^u_{n+1}U(t_{n+1},t)$. Their product is the invertible
two-dimensional block $B_n=U(t_{n+1},t_n)|_{E_n^u}$, so
$\Pi^u(t)=J_tB_n^{-1}F_t$ and $\Pi^c(t)=I-\Pi^u(t)$.
This proves invariance and bounds without inverting a parabolic map
on the infinite-dimensional space. Consequently
\begin{equation}
 z(s)=\int_0^s U(s,t)\Pi^c(t)f(t)\,dt
       -\int_s^T U(s,t)\Pi^u(t)f(t)\,dt
 \label{short-bg:green}
\end{equation}
satisfies the equation, has initial value in the fixed entrance pair,
has vanishing terminal growing coordinate, and obeys
\begin{equation}
 \sup_{[0,T]}\|z(s)\|_{X_m}
       \le C e^{a'T}\|f\|_{L^1(0,T;X_m)}.
 \label{short-bg:green-bound}
\end{equation}
Only a rank-two evolution was inverted. All constructions preserve
zero circulation and reality.

\subsection{Completion for arbitrary body forces}

Let $F\in H_p$, set $b_0=T_0F$, and solve for the correction in the $(r,z)$-components
\begin{equation}
 z_s-\nu\Delta z+\mathbb P_0(v\cdot\nabla z+z\cdot\nabla v)
     =g:=-\mathbb P_0(v\cdot\nabla b_0+b_0\cdot\nabla v).
 \label{short-bg:body-eq}
\end{equation}
Use the fixed Green selector \eqref{short-bg:green} for the vorticity
$\operatorname{curl}z$: its initial value belongs to the fixed entrance
pair and its terminal growing coordinate vanishes. No condition is
imposed on $b_0(T)$ or on a velocity observation. The train coefficient
and source remain present at every age. The global free graph gives
\[
 \|b_0\|_{H_p}+\|\sqrt{\nu s}Db_0\|_{H_p}
                    +\|\nu sD^2b_0\|_{H_p}\le C\|F\|_{H_p}.
\]
Put $b=13/8$, so $b-p=-9/8$. Spatial H\"older with $22/3$,
followed by one derivative, bounds both $g$ in $H^1$ and
$\operatorname{curl}g$ in $X_m$ on a finite slow-age band by
$\mathcal P_\nu(\|D^2b_0\|_p+\|Db_0\|_p+\|b_0\|_p)$.
Thus
\begin{equation}
 \|g\|_{L^p((0,T),s^bds;H^1)}
 +\|\operatorname{curl}g\|_{L^p((0,T),s^bds;X_m)}
 \le\mathcal P_\nu\|F\|_{H_p}.
 \label{short-bg:source}
\end{equation}
The circulation of the curl source is zero. The elementary estimates
\begin{align}
 \|f\|_{L^1(0,T)}&\le14^{7/11}T^{1/22}
                        \|f\|_{L^p(s^{13/8}ds)},\\
 \left\|\int_0^s f(t)\,dt\right\|_{L^p(s^{-9/8}ds)}
          &\le22\|f\|_{L^p(s^{13/8}ds)},\\
 \left\|\int_0^s(s-t)^{-1/2}f(t)\,dt\right\|_{L^p(s^{13/8}ds)}
          &\le C\sqrt T\|f\|_{L^p(s^{13/8}ds)}
 \label{short-bg:hardy}
\end{align}
follow respectively by H\"older and by $t=su$, Minkowski, and the
integrable factors $u^{-21/22}$ and
$(1-u)^{-1/2}u^{-21/22}$. These apply to norms of Banach-valued
functions.

The zero-circulation cylinder Biot--Savart map sends $X_m$ to $H^1$
with polynomial $c$ cost. On nonzero horizontal modes use the gap
$2\pi/c$; on the mean mode use
$|\widehat\omega(0,\xi)|\le|\xi|\int|\eta\bar\omega(\eta)|d\eta$
before dividing by $\xi$. Thus \eqref{short-bg:green-bound} and
\eqref{short-bg:source} give
\begin{equation}
 \sup_{[0,T]}\|z(s)\|_{H^1}+|\alpha|
 \le\mathcal P_\nu e^{a'T}\|F\|_{H_p},
 \label{short-bg:early}
\end{equation}
where $\alpha\in\mathbb R^2$ is the entrance coefficient of the
Green-selected correction. It is determined linearly by $F$; no force
trace is used.

For the global continuation, the train estimates imply, with normalized
spatial measure and $r=22/3$,
\[
 \|D^jv(s)\|_r\le\mathcal P_\nu s^{-(41+22j)/44},\quad
 \|v(s)\|_\infty\le\mathcal P_\nu s^{-1},\quad
 \|Dv(s)\|_\infty\le\mathcal P_\nu s^{-3/2},\qquad s\ge1.
\]
Write $B_j(s)=(\nu s)^{j/2}\|D^jb_0(s)\|_p$. H\"older gives
\begin{equation}
 \|g(s)\|_{H^1}\le\mathcal P_\nu
 [\nu^{-1}s^{-85/44}B_2(s)
       +s^{-63/44}(\nu^{-1/2}B_1(s)+B_0(s))].
 \label{short-bg:tail-source}
\end{equation}
Each $B_j$ has $L^p(s^{-9/8}ds)$ norm at most $C\|F\|_{H_p}$.
For a factor $s^{-k}B_j$, weighted H\"older on $(T,\infty)$ gives
$T^{23/22-k}$; hence the exponents in its $L^1$ bound are
$-39/44$ and $-17/44$, respectively. The same expression belongs to
$L^p(s^{-9/8}ds;H^1)$. Differentiated velocity energy and
\eqref{short-bg:train-late} therefore give
\begin{equation}
 M_z:=\sup_{s\ge0}\|z(s)\|_{H^1}
 \le\mathcal P_\nu
 \exp\!\left({a'\Sigma+C/(1+\Sigma)\over\nu}+C\nu^3\right)
 \|F\|_{H_p}.
 \label{short-bg:global-energy}
\end{equation}
Local smoothing introduces no further exponential. To check this,
differentiate \eqref{short-bg:body-eq} homogeneously twice. Incompressibility,
orthogonality of $\mathbb P_0$, and \eqref{short-bg:train-uniform} give
$E_j'+\nu E_{j+1}\le C\sum_{i\le j}E_i$, $j\le2$.
Multiplication by successive powers of $\nu(t-s)$ and summation
gives $H^1\to H^2$ norm $C[\nu(t-s)]^{-1/2}$ for $t-s\le1$.
On $0<s<2$, Duhamel and \eqref{short-bg:hardy} give the weighted
$L^p(s^{13/8}ds;H^2)$ estimate. For $s\ge2$ use Duhamel only on
$[s-1,s]$; weights at $s$ and $s-r$, $0<r<1$, are comparable.
Consequently
\[
 \|z\|_{L^p((T,\infty),s^{-9/8}ds;H^2)}
 \le C\nu^{-1/2}\bigl[T^{-1/22}M_z
             +\|g\|_{L^p((T-1,\infty),s^{-9/8}ds;H^1)}\bigr].
\]
Since $s\|v(s)\|_\infty+s\|Dv(s)\|_\infty\le\mathcal P_\nu$
for $s\ge1$, these estimates bound $sB_vz$ in $H_p$ with the right
side of \eqref{short-bg:global-energy}. Here
$B_vz=v\cdot\nabla z+z\cdot\nabla v$.

Let $E(s)\alpha$ be the fixed inlet stream with an age cutoff equal
to one near zero and zero before the terminal time $T$. Its homogeneous
solution differs from $E\alpha$ by $O(s)|\alpha|$ in the velocity
state. The first Hardy estimate controls the forced difference, so
$z-E\alpha\in H_p$ at zero; the global $H^1$ bound and
$\int_T^\infty s^{-9/8}ds=8T^{-1/8}$ give the tail.
Thus the exact global identity is
\begin{equation}
 b_0+z=E\alpha+T_0 f,\qquad
 f=F-sB_v(b_0+z)-s(\partial_s-\nu\Delta)E\alpha.
 \label{short-bg:flat-reconstruction}
\end{equation}
It follows from the differential equation and free uniqueness; each
term defining $f$ belongs to $H_p$. For the remaining free interaction
term this is explicit: on $s<T$,
$\|sB_vb_0\|_{H_p}\le\mathcal P_\nu(\nu^{-1/2}\sqrt T+T)\|F\|_{H_p}$;
on $s>T$, the bounds $v=O(s^{-1})$, $Dv=O(s^{-3/2})$ and the free
weighted gradient give
$\|sB_vb_0\|_{H_p}\le\mathcal P_\nu T^{-1/2}(\nu^{-1/2}+1)\|F\|_{H_p}$.
Approximation in
$L^p(s^{13/8}ds;X_m)$ and the difference estimates above justify all
identities for the original arbitrary force. Choosing $\Sigma$ large
and then $a'$ small proves, for every fixed $m>0$,
\begin{equation}
 \|f\|_{H_p}+|\alpha|\le\mathcal P_\nu e^{m/\nu}
                                \|F\|_{H_p}.
 \label{short-bg:flat-inverse}
\end{equation}

\subsection{Exact geometry and the nonlinear correction}

The exact stream lift $E_\epsilon\alpha$ is affine: its nonzero
inlet is not in $Y_p$. Its complete equation source is in $Y_p$.
The Gaussian angular-two eigenstream is a harmonic quadrupole plus
a Gaussian remainder outside a fixed core. Periodizing its harmonic
part uses only
\[
 \sum_{n\in\mathbb Z}(q+nc+i\eta)^{-2}
       =(\pi/c)^2\csc^2(\pi(q+i\eta)/c).
\]
Its apparent poles are replaced by the smooth eigenstream core. Thus
the actual periodized stream and every fixed derivative have
exponential transverse decay, with polynomial $c$ costs. This proves
the inlet bounds needed below without any estimate on a dual tail.
Indeed write $K=2s\partial_s+\eta\partial_\eta$, and for one inlet
stream $\psi$ write its normalized lift as
\[
 e_\epsilon=e_0+\epsilon e_1,\qquad
 e_0=\chi(-\psi_\eta,\psi_q),\qquad
 e_1=(0,(\chi-2s\chi')\psi-\chi\eta\psi_\eta).
\]
The exact vector diffusion is
$L_\epsilon=\Delta-2\epsilon(K+1)\partial_q
+\epsilon^2[(K+1)^2-\operatorname{diag}(1,0)]$.
Substitution in the complete unprojected source column
\[
 \mathcal C_\epsilon\alpha
 =s[(\partial_s-\nu L_\epsilon)e_\epsilon+B_\epsilon e_\epsilon]\alpha
\]
gives $\mathcal C_\epsilon-\mathcal C_0=\epsilon s r_\epsilon$,
where $r_\epsilon$ is smooth at zero, compact in upper age, and has
polynomial weighted bounds. Hence
\begin{equation}
 \|\mathcal C_\epsilon-\mathcal C_0\|_{H_p}\le\epsilon\mathcal P_\nu,
 \qquad
 \|A_0E_\epsilon\alpha+\mathbb PB_{V_0}E_\epsilon\alpha\|_{W^{1,Y_p}}
 \le\mathcal P_\nu R^{2/p}|\alpha|.
 \label{short-bg:affine}
\end{equation}
The integrability exponent is $-9/8+p=13/8>-1$. The quadratic
inlet source has the same bound with $|\alpha|^2$, and
$\|E_\epsilon\alpha\|_{W^{j,\infty}}\le\mathcal P_{\nu,j}|\alpha|$.
These assertions compare sources, never the nonintegrable inlet states.

The global flat completion now permits one all-age perturbation. Put
$X=H_p$, $K_\epsilon=sB_\epsilon T_\epsilon$,
$K_0=sB_vT_0$. The flat train is retained at every age. Give
$X\times\mathbb R^2$ the norm $\|(f,\alpha)\|=\|f\|_X+|\alpha|$.
The equation row is
\[
 \mathfrak A_0:X\times\mathbb R^2\longrightarrow X,\qquad
 \mathfrak A_0(f,\alpha)=(I+K_0)f+\mathcal C_0\alpha.
\]
The all-age reconstruction \eqref{short-bg:flat-reconstruction}
and the complete inlet source defining $\mathcal C_0$ give
$\mathfrak A_0(f,\alpha)=F$. The Green selector, Biot--Savart map
and reconstruction are linear, so \eqref{short-bg:flat-inverse}
constructs a bounded linear right inverse
\[
 R_0:X\longrightarrow X\times\mathbb R^2,\qquad
 \mathfrak A_0R_0=I_X,\qquad
 \|R_0\|\le M=\mathcal P_\nu e^{m/\nu}.
\]
We compare operator outputs while inputs remain arbitrary global
elements of $H_p$.

Here the all-age mean expansion gives a useful stronger coefficient
bound than a fixed small late constant:
\begin{equation}
 \sup_{s>S}(|V_0|+|DV_0|)
       \le\mathcal P_\nu(S^{-1/2}+\epsilon).
 \label{short-bg:strong-late-coeff}
\end{equation}
To check the Cartesian gradient, write $b=(v_1,v_2+\epsilon H)$.
Its three derivative types are $s\partial_\eta b$,
$s[\partial_q-\epsilon(K+1)]b$, and the frame term
$\epsilon s v_1$. The mean stream is
$S_0=a_0G_{\nu s}+a_2\partial_\eta^2G_{\nu s}
+O_{\rm weighted}(\mathcal P_\nu s^{-5/2})$.
Thus $v_1=-\partial_\eta S_0$ and
$H=2a_0G_{\nu s}+4a_2\partial_\eta^2G_{\nu s}+\cdots$.
The potentially nondecaying radial derivative cancels because
$(K+1)G_{\nu s}=0$, giving
$(K+1)H=-8a_2\partial_\eta^2G_{\nu s}
+O_{\rm weighted}(\mathcal P_\nu s^{-5/2})$.
Every fluctuation derivative is exponentially small. These identities
prove \eqref{short-bg:strong-late-coeff}, including the full frame
term; the axial mean velocity itself has the genuine $O(\epsilon)$
floor. The physical free graph and the exact norm identity then give
\[
 \|1_{s>S}K_\epsilon\|_{X\to X}
 \le\mathcal P_\nu(S^{-1/2}+\epsilon).
\]
For the flat operator, $v=O(\mathcal P_\nu/s)$ and
$Dv=O(\mathcal P_\nu/s^{3/2})$, paired with the free
$\sqrt{s}Db_0$ graph bound, give
$\|1_{s>S}K_0\|\le\mathcal P_\nu S^{-1/2}$.
At the young end, smooth inlet coefficients give
$\|1_{s<\sigma}K_\epsilon\|+\|1_{s<\sigma}K_0\|
\le\mathcal P_\nu\sqrt\sigma$ by the same free graph.

Take $\sigma=\epsilon^{a}$ and $S=\epsilon^{-b}$, where $a,b>0$
will be small fixed numbers. In \eqref{sct:transverse}, choose
$l=C_D(c/\nu+\sqrt S+1)$. Both negative exponents dominate any
prescribed multiple of $\nu^{-1}$, while $P_S\le\mathcal P_\nu
e^{kb d_+/\nu}$ for fixed $k$. Increasing $C_D$ makes all required
coefficient tails at most $e^{-D/\nu}$; fixed weights and $K$
derivatives change only the polynomial cost. The prefactor in the
localized comparison is polynomial in this width, $S$ and
$\sigma^{-1}$. With $\theta=2/539$ from
Proposition~\ref{bg:lp-free-comparison}, for some fixed $N<\infty$,
\begin{equation}
 \|K_\epsilon-K_0\|_{X\to X}
 +\|\mathcal C_\epsilon-\mathcal C_0\|
 \le\mathcal P_\nu
 [\epsilon^{a/2}+\epsilon^{b/2}+\epsilon
       +\epsilon^{\theta-N(a+b)}+e^{-D/\nu}]
 =:\Delta_\nu.
 \label{short-bg:global-operator-comparison}
\end{equation}
The affine column is supported
in a fixed polynomial age interval and has the already proved
$O(\epsilon)$ source difference. No equation or source is cut off.
Choose $a,b$ so $\theta-N(a+b)>0$, and put
$\zeta=\min\{a/2,b/2,1,\theta-N(a+b)\}>0$.
Choose $m<d_-\zeta/2$ and $D>2m$. Then $M\Delta_\nu<1/2$ for
small $\nu$. The exact equation row
\[
 \mathfrak A_\epsilon:X\times\mathbb R^2\longrightarrow X,\qquad
 \mathfrak A_\epsilon(f,\alpha)
       =(I+K_\epsilon)f+\mathcal C_\epsilon\alpha
\]
has the bounded linear right inverse
\[
 R_\epsilon=R_0[I+(\mathfrak A_\epsilon-\mathfrak A_0)R_0]^{-1},
 \qquad \|R_\epsilon\|\le2M.
\]
Indeed $(\mathfrak A_\epsilon-\mathfrak A_0)R_0:X\to X$
has norm at most $M\Delta_\nu<1/2$, and multiplication verifies
$\mathfrak A_\epsilon R_\epsilon=I_X$.
This compares all global inputs, including those supported arbitrarily
late; it does not assume that the exact inverse is causal in age.
Returning to physical normalization proves
\begin{equation}
 \|w\|_{X_p}+R^{2/p}|\alpha|
 \le\mathcal P_\nu e^{m/\nu}
                \|f\|_{Y_p}
 \label{short-bg:exact-inverse}
\end{equation}
for
$A_0w+\mathbb PB_{V_0}w+(A_0E_\epsilon+\mathbb PB_{V_0}E_\epsilon)\alpha
=\mathbb Pf$. Here the physical $w$ is the free response to the
first coordinate selected by $R_\epsilon$, and the free graph supplies
its $X_p$ norm. The same normalized measure multiplies
source, state, and finite coordinate, so no additional power of $R$
is lost.

Set $A=R^{2/p}\alpha$ and $\|(w,A)\|_Z=\|w\|_{X_p}+|A|$.
By Lemmas~\ref{bg:common-body-graph}--\ref{bg:graph-embeddings} and
\eqref{short-bg:affine}, the quadratic
terms in $w+E_\epsilon\alpha$ and their polarized differences are
bounded by $\mathcal P_\nu\|(w,A)\|_Z^2$ in $Y_p$. Apply the fixed
linear right-inverse selection to minus these terms and minus $F_R$.
With $\delta=\mathcal P_\nu\epsilon^{3/11}$ and
$M=\mathcal P_\nu e^{m/\nu}$, this is a contraction on the radius
$2M\delta$ ball if $4\mathcal P_\nu M^2\delta<1$. This holds after
decreasing $m$ so $2m<3d_-/11$. We obtain an exact smooth unforced
profile
\begin{equation}
 U=V_0+E_\epsilon\alpha+w,\qquad
 \|w\|_{X_p}+R^{8/11}|\alpha|
       \le\mathcal P_\nu e^{m/\nu}\epsilon^{3/11},\qquad
 |\alpha|\le\mathcal P_\nu e^{m/\nu}\epsilon.
 \label{short-bg:result}
\end{equation}
The projected equation is
$A_0U+\mathbb P(U\cdot\nabla U)=0$ with zero residual. Reality,
no-swirl symmetry, and periodicity are preserved at each iteration.

No second unstable inverse is needed for spatial regularity. Write
$V=V_0+E_\epsilon\alpha$ and
$A_0w=f-\mathbb P(V\cdot\nabla w+w\cdot\nabla V+w\cdot\nabla w)$.
The complete source has $W^{1,Y_p}$ norm bounded by the right side
of \eqref{short-bg:result}, up to a polynomial. Since
$\partial_jA_0w=(A_0-1/2)\partial_jw$ and
$\|Dw\|_{Y_{2p}}^2\le C\|w\|_\infty\|D^2w\|_{Y_p}$,
every differentiated term is in $Y_p$. Also $Dw\in Y_p$ is already
known. Rewrite the differentiated equation as
$A_0D_jw=f_j+\tfrac12D_jw$, with the entire right side in $Y_p$,
and apply the maximal graph at shift zero. No inverse of
$A_0-\tfrac12$ is used. Thus $Dw\in\widetilde X_p$, and
\begin{equation}
 \sup_\tau\|w(\tau)\|_p+
 \|w\|_{W^{1,\infty}}+\|E_\epsilon\alpha\|_{W^{1,\infty}}
 \le\mathcal P_\nu e^{m/\nu}\epsilon^{3/11}.
 \label{short-bg:strong}
\end{equation}
For any fixed derivative order, the train and inlet complete sources
belong to the corresponding $W^{k,Y_p}$. Explicitly, in the
large-radius region their unprojected derivatives are bounded by
$C_{k,N}\varrho^{-3-k}(1+|y_3|/\varrho)^{-N}$, whose $p$th integral
reduces to $C\int_R^\infty\varrho^{2-p(3+k)}d\varrho<\infty$.
On $\varrho<R$, take the all-age heat expansion to a sufficiently
high finite order: its even radial powers extend smoothly through the
axis and its weighted remainder gives arbitrary axial decay. The
inlet cutoff vanishes near the axis. These estimates are uniform in
phase; the bounded Cartesian Leray map preserves them in $W^{k,p}$.
Inductively, if derivatives
through order $k$ are in the graph, all factors through order $k$ are
bounded and derivatives through order $k+2$ are in $Y_p$. In each
order-$(k+1)$ derivative of $w\cdot\nabla w$, one factor has order at
most $k$. Put it in $L^\infty$ and the other in $Y_p$, then apply the
maximal graph to the differentiated equation. This proves all spatial
orders. On positive physical-time compact intervals the resulting
$C_\tau W^{k,4}$ regularity and the projected equation give all time
derivatives, so $U$ is jointly smooth through the axis.

\subsection{The datum, heat remainder, and pressure}

Let $\psi_{\rm in}$ be the raw inlet plus its single entrance
adjustment. Define, with $\varrho=|x'|$,
\[
 q=R\log(\varrho/R),\quad\eta=Rz/\varrho,\qquad
 a_r={R\over\varrho}(-\psi_{{\rm in},\eta}),\quad
 a_z={R\over\varrho}
 [\psi_{{\rm in},q}+\epsilon(\psi_{\rm in}
                          -\eta\psi_{{\rm in},\eta})].
\]
The stream is periodic in $q$ and has arbitrary fixed transverse
polynomial decay. Differentiation and this decay give
$|D^ja(x)|\le C_j|x|^{-1-j}$, including smooth extension at each
nonzero axis point. The omitted divergence flux through a radius-$r$
sphere is $O(r)$, so $a$ is distributionally divergence-free on all of
$\mathbb R^3$. Thus $a\in L^{3,\infty}\cap L^2_{\rm loc}$.
The shift $q\mapsto q+c$ proves its DSS identity. The raw inlet is
nonzero on a fixed core and \eqref{short-bg:result} makes its finite
adjustment exponentially small there, so $a\ne0$.

Put $a_\tau(y)=e^{\tau/2}a(e^{\tau/2}y)$ and
$H_a(\tau)=e^{\nu\Delta}a_\tau$. Taylor expansion against the Gaussian,
whose first moment vanishes, gives
\[
 |D^j(e^{\nu\Delta}a_\tau-a_\tau)(y)|
       \le C_{j,\nu}|y|^{-3-j},\qquad |y|\ge1.
\]
Use the expansion only in a ball of radius $|y|/4$ about $y$;
differentiate the Gaussian on its complement, which is exponentially
small despite the singularity of $a$ at zero. The $O(s)$ inlet
expansion gives the same inverse-cubic bound for $V-a_\tau$ in the
large-radius region, with arbitrary transverse decay. The all-age
train expansion gives arbitrary axial decay in the remaining bounded
radial region. Therefore
\begin{equation}
 V-H_a\in C_{\rm per}H^k\cap X_p
 \quad\hbox{for every fixed required }k.
 \label{short-bg:raw-heat}
\end{equation}
The graph membership includes the mixed phase/dilation derivative;
it follows from the same source expansion. In particular this
conclusion is stronger than an $X_p$ estimate alone.

The inverse-distance bounds and the all-age expansion give
$D^jV\in C_{\rm per}L^4$. For $w$, interpolation between its
$C_{\rm per}L^p$ derivatives and their uniform supremum bounds gives
$D^jw\in C_{\rm per}L^4$. Therefore
$D^jU\in C_{\rm per}L^4$ for $j\le2$, and
$\mathbb T=U\otimes U\in C_{\rm per}H^2$. The field
$Q=U-H_a\in X_p$ solves $A_0Q=-\mathbb P\operatorname{div}\mathbb T$.
The exact free semigroup is
\[
 S_0(t)f(y)=e^{t/2}[e^{\nu(e^t-1)\Delta}f](e^{t/2}y).
\]
Free uniqueness identifies $Q$ with
$-\int_0^\infty S_0(t)\mathbb P\operatorname{div}\mathbb T(\tau-t)dt$.
For $j\le2$, move derivatives to $\mathbb T$ when $t\le1$ and to
the heat kernel when $t\ge1$. The resulting integrable majorants are
\[
 \|D^jS_0(t)\mathbb P\operatorname{div}\mathbb T\|_2
 \le C_{j,\nu}
 \begin{cases}t^{-1/2}\|\mathbb T\|_{H^j},&0<t\le1,\\
 e^{-3t/4}\|\mathbb T\|_2,&t\ge1.
 \end{cases}
\]
Thus $Q\in C_{\rm per}H^2$. Also
$(\partial_\tau-y\cdot\nabla/2)Q
=\nu\Delta Q+Q/2-\mathbb P\operatorname{div}\mathbb T
\in C_{\rm per}L^2$. Scaling back gives
\[
 \|u_\nu(t)-e^{\nu t\Delta}a\|_2
       =t^{1/4}\|Q(\log t)\|_2\le C t^{1/4}.
\]
A compact part of $a$ is in $L^2$; its complement is bounded and
uniformly continuous. The Gaussian approximate identity therefore
shows $e^{\nu t\Delta}a\to a$ locally in $L^2$, and hence the same
is true for $u_\nu$. Finally
$p_U=\sum_{i,j}R_iR_j(U_iU_j)\in C_{\rm per}L^2$ recovers the
discarded gradient part of the exact projected equation.
